\subsection*{3.2 Pruning a depth four circuit}
As mentioned before, we will prune the circuit $\mathtt C$ computing $f_n$ in two steps--first we will restrict the bottom support of $\mathtt C$ and then we will get rid of all heavy gates in it.
\subsection*{3.2.1 Step1--Restricting the bottom support of $\mathtt C$}
If the number of monomials in $\mathtt C$ is more than $\lfloor m^2d_x/(\ln m)^5\rfloor$, there is nothing to prove. Otherwise, we show that we can get rid of all monomials with support more than $\tau=\lfloor20\ln m\rfloor$.
\begin{authorclaim}{Lemma10}
Let the number of monomials in $\mathtt C$ be at most $\lfloor m^2d_x/(\ln m)^5\rfloor$. Then, for sufficiently large $m$, there exists $M_1\subseteq[3m]$, $|M_1|=m$ such that all monomials in $\mathtt C_1$ obtained from $\mathtt C$ by setting variables $\mathbf x_{M_1}$ and $\mathbf y_{M_1}$ to 0 have support at most $\tau$.
\end{authorclaim}
\begin{proof}
We first present a greedy procedure to remove all monomials with support more than $\tau$ and then argue that it sets $m$ variables each from $\mathbf x$ and $\mathbf y$ to 0. In each iteration the procedure picks a pair of variables that appears in a large number of monomials with support more than $\tau$ and set them to 0.
\begin{authorclaim}{Procedure1:Restriction procedure}
\begin{enumerate}
\item $M_1\leftarrow\varnothing$, $\mathtt C_1\leftarrow\mathtt C$, $H:=$ set of all monomials of $\mathtt C_1$ with support more than $\tau$.
\item For $j\in[3m]$, $e(j):=$ number of monomials in $H$ containing $x_j$ or $y_j$.
\item while $H\ne\varnothing$ do
\item Pick $j'\in[3m]\setminus M_1$ such that $e(j')\ge e(j)$ for all $j\in[3m]$. Set $x_{j'}=0$ and $y_{j'}=0$. Update $M_1\leftarrow M_1\cup\{j'\}$, $\mathtt C_1\leftarrow$ circuit obtained from $\mathtt C_1$ by setting $x_{j'}$ and $y_{j'}$ to 0, $H\leftarrow$ set of all monomials of $\mathtt C_1$ with support more than $\tau$, and $e(j)\leftarrow$ number of monomials in $H$ containing $x_j$ or $y_j$.
\item end while
\end{enumerate}
\end{authorclaim}
It is clear that the bottom support of $\mathtt C_1$ obtained after the termination of the procedure is at most $\tau$. Also, since we are only setting variables to 0, it trivially follows that the procedure does not increase the number of gates nor does it increase the fan-in of any gate in the circuit. Claim11 (proved in AppendixB.2) implies that the procedure terminates in at most $m$ iterations. If it terminates before $m$ iterations, we arbitrarily add an appropriate number of $j\in[3m]$ to $M_1$ so that $|M_1|=m$ and set $x_j$ and $y_j$ to 0 for all such $j$.
\end{proof}
\begin{authorclaim}{Claim11}
Procedure1 terminates in at most $m$ iterations.
\end{authorclaim}
\paragraph{Remark}
Procedure1 looks similar to an approximation algorithm for the Set Cover problem [102] (Section1.2,page6). This is because the problem of removing monomials with support more than $\tau$ can be formulated as an instance of Set Cover with the universe being all such monomials and with a set corresponding to each $j\in[3m]$. For a $j\in[3m]$, the corresponding set will contain all monomials with support more than $\tau$ in which at least one of $x_j$ and $y_j$ appears.
\subsection*{3.2.2 Step2--Pruning the heavy gates from $\mathtt C_1$}
A sparse polynomial $Q$ in $\mathtt C_1$ is said to be light if its fan-in is at most $m/(\ln m)^2$, i.e., at most $m/(\ln m)^2$ many non-zero monomials are present in $Q$. If the sum of fan-ins of all the light sparse polynomials is $\Omega(m^2d_x/(\ln m)^5)$ then there is nothing to prove. So assume that the sum of fan-ins of all the light sparse polynomials is $O(m^2d_x/(\ln m)^5)$. Recall that a product term is called heavy if it has at least $w=\lfloor md_x/(\lambda_0(\ln m)^3)\rfloor$ many distinct sparse polynomials connected to it (where $\lambda_0$ is a large enough constant fixed in AppendixC). Observe that one of the following cases is true:
\begin{enumerate}
\item There is a heavy gate in $\mathtt C_1$, that is connected to at most $md_x/(2\lambda_0(\ln m)^3)$ light sparse polynomials.
\item Every heavy gate is connected to at least $md_x/(2\lambda_0(\ln m)^3)$ light sparse polynomials.
\end{enumerate}
Case1 clearly implies a lower bound of $\Omega(m^2d_x/(\ln m)^5)$. Else, we prove the following lemma.
\begin{authorclaim}{Lemma12}
Let $\mathtt C_1$ be the circuit (obtained from Lemma10) having at most $m$ heavy gates such that every heavy gate is connected to at least $md_x/(2\lambda_0(\ln m)^3)$ light sparse polynomials, and the sum of fan-ins of all the light sparse polynomials is at most $m^2d_x/(160\lambda_0(\ln m)^5)$. Then, there exist $M_2\subseteq[3m]\setminus M_1$, $|M_2|=m$ and $\alpha_l,\beta_l\in\mathbb F$ for $l\in M_2$, such that setting $x_l=\alpha_l$, $y_l=\beta_l$ for all $l\in M_2$ removes all heavy gates from $\mathtt C_1$.
\end{authorclaim}
\begin{proof}
We first present the pruning procedure and then argue its correctness. For any light sparse polynomial $Q_j$ in $\mathtt C_1$, let $b_j$ and $c_j$ be equal to the fan-in of $Q_j$ and the number of distinct heavy gates connected to $Q_j$ in $\mathtt C_1$ respectively. As $Q_j$ is a light sparse polynomial, $b_j\le m/(\ln m)^2$. In this procedure, while all the heavy gates do not disappear from $\mathtt C_1$, we pick a sparse polynomial whose ratio of the number of distinct heavy gates connected to it and fan-in is maximum and set that to zero by evaluating it to one of its roots in $\mathbb F$. This can be done as $\mathbb F$ is algebraically closed. The restricted support of the monomials in $\mathtt C_1$ ensure that at the end of this procedure, only a small fraction of the variables are set.
\begin{authorclaim}{Procedure2:Pruning heavy gates from $\mathtt C_1$}
\begin{enumerate}
\item Set $i=1$ and $M_2=\varnothing$. Let $s_1\le m$ be the number of heavy gates in $\mathtt C_1$. Choose a non-constant light sparse polynomial $Q_1$ from $\mathtt C_1$ such that the ratio $c_1/b_1$ is maximum and add the indices of the variables appearing in $Q_1$ to $M_2$. As $b_1\le m/(\ln m)^2$, we have $\tau b_1\le m$.
\item Make $Q_1$ equal to zero by setting at most $\tau b_1$ many variables appearing in $Q_1$ to field constants. By doing so, at least $c_1$ many heavy gates vanish from $\mathtt C_1$.
\item while $(\tau(b_1+\cdots+b_i)\le m)$ do
\item Increment $i$ by 1.
\item Let $s_i$ be the number of heavy gates in the current circuit $\mathtt C_1$ obtained after the $(i-1)$-th iteration. Clearly, $s_i\le s_{i-1}-c_{i-1}$. If $s_i=0$ then exit the loop.
\item Otherwise, choose a non-constant light sparse polynomial $Q_i$ from $\mathtt C_1$ having the maximum value of $c_i/b_i$ in $\mathtt C_1$ and add the indices of the variables appearing in $Q_i$ to $M_2$.
\item Make $Q_i$ equal to zero by setting at most $\tau b_i$ many variables appearing in $Q_i$ to field constants. By doing so, at least $c_i$ many heavy gates vanish from $\mathtt C_1$.
\item end while
\end{enumerate}
\end{authorclaim}
\begin{authorclaim}{Claim13}
Let $\overline M_1=[3m]\setminus M_1$. Procedure2 sets at most $m$ many variables in $\mathbf x_{\overline M_1}\cup\mathbf y_{\overline M_1}$ to field constants and removes all the heavy gates from $\mathtt C_1$.
\end{authorclaim}
The above claim is proved in AppendixB.3. The claim implies that $|M_2|\le m$. If $|M_2|<m$, add appropriate number of elements from $[3m]\setminus M_1\cup M_2$ to $M_2$ arbitrarily so that $|M_2|=m$. For every $l\in M_2$, if $x_l$ or $y_l$ is not set to a field constant then set $x_l=0$ or $y_l=0$ respectively. Clearly, we end up setting exactly $m$ variables each from $\mathbf x_{\overline M_1}$ and $\mathbf y_{\overline M_1}$.
\end{proof}
\paragraph{Remark}
Procedure2 resembles an approximation algorithm for the Weighted Set Cover problem [99] (Section2.1,page16). This is no coincidence as the problem of removing heavy gates can be formulated as an instance of Weighted Set Cover with the universe being all heavy gates and with a set corresponding to every sparse polynomial $Q$. The set corresponding to $Q$ contains all heavy gates connected to $Q$ and has a cost equal to the number of monomials feeding into $Q$.
